\documentclass[11pt]{article}

\usepackage[margin=1in]{geometry}
\usepackage{amsmath, amssymb}
\usepackage[T1]{fontenc}
\usepackage{lmodern}
\usepackage{microtype}
\usepackage{graphicx}
\usepackage{booktabs}
\usepackage{array}
\usepackage{float}
\usepackage{xcolor}
\usepackage{fancyhdr}
\usepackage[font=small,labelfont=bf,skip=5pt]{caption}
\usepackage{listings}
\lstdefinestyle{pycode}{
  language=Python, basicstyle=\ttfamily\scriptsize,
  keywordstyle=\color{psink}\bfseries, commentstyle=\color{psaccent}\itshape,
  stringstyle=\color{psink}, numbers=left, numberstyle=\tiny\color{psrule},
  numbersep=7pt, breaklines=true, breakatwhitespace=true, showstringspaces=false,
  frame=none, xleftmargin=20pt, aboveskip=6pt, belowskip=6pt, columns=fullflexible,
  keepspaces=true, upquote=true,
}
\lstset{style=pycode}

\definecolor{psink}{RGB}{23,54,93}
\definecolor{psaccent}{RGB}{192,80,77}
\definecolor{psrule}{RGB}{200,209,222}

\usepackage[colorlinks=true,allcolors=psink]{hyperref}

\graphicspath{{./}}
\emergencystretch=3em

\setlength{\textfloatsep}{10pt plus 2pt minus 2pt}
\setlength{\intextsep}{10pt plus 2pt minus 2pt}
\setlength{\abovecaptionskip}{4pt}
\setlength{\belowcaptionskip}{0pt}

\makeatletter
\renewcommand\section{\@startsection{section}{1}{0pt}%
  {-2.6ex \@plus -1ex \@minus -.2ex}{1.3ex \@plus .2ex}%
  {\normalfont\scshape\large\color{psink}}}   % lmr has no bold small caps
\renewcommand\subsection{\@startsection{subsection}{2}{0pt}%
  {-2.2ex \@plus -1ex \@minus -.2ex}{1.0ex \@plus .2ex}%
  {\normalfont\bfseries\normalsize\color{psink}}}
\renewcommand\subsubsection{\@startsection{subsubsection}{3}{0pt}%
  {-1.9ex \@plus -.8ex \@minus -.2ex}{0.8ex \@plus .2ex}%
  {\normalfont\bfseries\itshape\normalsize\color{psink}}}
\renewcommand\subparagraph{\@startsection{subparagraph}{5}{0pt}%
  {-1.6ex \@plus -.6ex \@minus -.2ex}{0.5ex \@plus .1ex}%
  {\normalfont\itshape\normalsize\color{psink}}}
\renewcommand\paragraph{\@startsection{paragraph}{4}{0pt}%
  {1.7ex \@plus 1ex \@minus .2ex}{-0.9em}%
  {\normalfont\bfseries\normalsize\color{psink}}}
\makeatother

\newenvironment{tight}
  {\begin{itemize}\setlength{\itemsep}{3pt}\setlength{\parskip}{0pt}%
   \setlength{\topsep}{3pt}\setlength{\leftmargini}{1.4em}}
  {\end{itemize}}

\setcounter{secnumdepth}{0}   % headings carry the assignment's own labels
\setlength{\headheight}{14pt}
\pagestyle{fancy}
\fancyhf{}
\fancyhead[L]{\footnotesize\scshape\color{psink} MGT 595 \textemdash\ Problem Set 5}
\fancyhead[R]{\footnotesize\color{psink}\thepage}
\renewcommand{\headrulewidth}{0.4pt}
\renewcommand{\headrule}{\hbox to\headwidth{%
  \color{psrule}\leaders\hrule height \headrulewidth\hfill}}

\begin{document}

\thispagestyle{empty}
\begin{flushleft}
{\bfseries\LARGE\color{psink} Quantitative Investing}\\[3pt]
{\scshape\Large\color{psink} Problem Set 5 \textemdash\ Part I: Business Cycle Variation}\\[10pt]
{\color{psrule}\rule{\textwidth}{1.2pt}}\\[5pt]
{\small Ashley Wen\,(xw564)\quad\textbar\quad Emily Wu\,(esw63)%
\quad\textbar\quad Grace Gu\,(xg325)\hfill \today}
\end{flushleft}

\vspace{1.4em}

% =====================================================================
\section{Data and Conventions}
% =====================================================================

\begin{tight}
  \item \textbf{Sample.} Monthly data from 1926-07 to 2026-06 (UMD and the momentum
  portfolios start 1927-01). The NBER dummy marks 201 recession months across 16 recessions.
  \item \textbf{Returns.} All figures are percent per month. Portfolio returns are in
  excess of the T-bill; RMRF is already excess, and SMB, HML and UMD are zero-cost
  long--short returns, so they are used as given.
  \item \textbf{Inference.} Residuals cluster in recession episodes, so $t$-statistics
  use Newey--West standard errors with 12 lags.
  \item \textbf{Shared code.} All parts load data through \texttt{common/data.py} and run
  Fama--MacBeth through \texttt{common/fmb.py}, so lags, samples and beta definitions
  are identical across Questions (a)--(h).
\end{tight}

% =====================================================================
\section{Question (a)}
% =====================================================================

We estimate $R_t = a + b\,D_t + \varepsilon_t$, where $D_t = 1$ in NBER recession months,
so $a$ is the non-recession mean and $a+b$ the recession mean.

\begin{table}[H]\centering\small
\caption{Factor returns in and out of recessions (\% per month)}
\begin{tabular}{lcccc}
\toprule
 & Non-recession $a$ & Recession $a+b$ & Difference $b$ & $t(b)$ \\
\midrule
RMRF & 0.93 & $-0.47$ & $-1.40$ & $-2.50$ \\
SMB  & 0.19 & 0.03    & $-0.16$ & $-0.82$ \\
HML  & 0.40 & 0.15    & $-0.25$ & $-0.75$ \\
UMD  & 0.71 & 0.22    & $-0.49$ & $-1.00$ \\
\bottomrule
\end{tabular}
\par\smallskip\footnotesize Newey--West (12 lags) $t$-statistics. UMD starts 1927-01.
\end{table}

\begin{tight}
  \item \textbf{What does this tell you?} Only the market is reliably worse in
  recessions; SMB, HML and UMD give up some return but none of the drops is
  distinguishable from zero. All three factors still earn a positive average return in
  recession months, while the market loses money. So the factor premia are not
  mainly a reward for losing money in recessions, at least not in the way the market premium is.
  \item \textbf{Based on these results, do you believe SMB or HML is more consistent
  with a risk-based explanation?} SMB is, but only weakly. Its premium all but
  disappears in recessions, so small stocks fail to pay off exactly when times are
  bad, which is what a risk story needs. HML instead keeps a positive return in
  recessions, in line with Lakonishok, Shleifer and Vishny's (1994) finding that value does not
  underperform in bad states. With 201 recession months neither difference is
  significant, so we read this as a lean rather than a verdict.
\end{tight}

% =====================================================================
\section{Question (b)}
% =====================================================================

\begin{table}[H]\centering\small
\caption{Corner portfolio excess returns in and out of recessions (\% per month)}
\begin{tabular}{lcccc}
\toprule
 & Non-recession $a$ & Recession $a+b$ & Difference $b$ & $t(b)$ \\
\midrule
\multicolumn{5}{l}{\textit{Size--BE/ME corners}} \\
Small Growth & 0.83 & $-0.63$ & $-1.46$ & $-1.30$ \\
Small Value  & 1.69 & $-0.38$ & $-2.07$ & $-2.62$ \\
Big Growth   & 0.90 & $-0.34$ & $-1.25$ & $-2.21$ \\
Big Value    & 1.24 & $-0.12$ & $-1.35$ & $-1.51$ \\
\addlinespace
\multicolumn{5}{l}{\textit{Size--momentum corners}} \\
Small Losers  & 0.80 & $-0.26$ & $-1.06$ & $-1.17$ \\
Small Winners & 1.84 & $-0.22$ & $-2.06$ & $-2.98$ \\
Big Losers    & 0.58 & $-0.83$ & $-1.42$ & $-1.66$ \\
Big Winners   & 1.24 & $-0.40$ & $-1.64$ & $-2.73$ \\
\addlinespace
\multicolumn{5}{l}{\textit{Within-size spreads}} \\
Small Value $-$ Growth   & 0.87 & 0.25 & $-0.62$ & $-1.03$ \\
Big Value $-$ Growth     & 0.34 & 0.23 & $-0.11$ & $-0.14$ \\
Small Winners $-$ Losers & 1.04 & 0.04 & $-1.00$ & $-1.29$ \\
Big Winners $-$ Losers   & 0.64 & 0.39 & $-0.25$ & $-0.40$ \\
\bottomrule
\end{tabular}
\par\smallskip\footnotesize Excess returns over the T-bill. Newey--West (12 lags)
$t$-statistics. Momentum portfolios start 1927-01.
\end{table}

\begin{tight}
  \item \textbf{Do you find stronger cyclical patterns for these extreme portfolios
  than for the factors?} Yes, but mostly because the corners are long-only portfolios that
  carry the market's recession losses. Every corner falls by one to two percent per
  month in recessions, and the high-return corners (Small Value and both
  Winner portfolios) fall the most. Once we net out the market by taking
  spreads within a size group, the pattern shrinks back to factor-sized, insignificant
  drops that stay positive in recessions. The extreme portfolios are more cyclical,
  but the value and momentum premia themselves are not.
\end{tight}

\end{document}
